\section{Object oriented programming}

  \ac{OOP} is a way of structuring software around \textit{objects}. An object combines data and behaviour, meaning that it stores information and also provides operations that can act on that information. Instead of treating a program as one long sequence of instructions, the program is divided into units that represent meaningful concepts in the problem domain \cite{ObjectorientedProgramming2026,JacobsonOOSE}.

  In practical terms this usually means that the programmer defines a \textit{class}, which acts as a template, and then creates \textit{objects} from that class. Figure~\ref{fig:oop_class_object} illustrates this idea using a simple melody example. The class describes what information a melody should contain and what operations are available, while the object is one concrete melody with specific values.

  \begin{figure}[htbp]
  \centering
  \resizebox{0.92\textwidth}{!}{%
  \begin{tikzpicture}[
    >=Latex,
    font=\small,
    node distance=12mm and 14mm,
    classbox/.style={
      draw,
      rounded corners=3pt,
      minimum width=58mm,
      align=left,
      fill=blue!6
    },
    objectbox/.style={
      draw,
      rounded corners=3pt,
      minimum width=58mm,
      align=left,
      fill=green!7
    },
    arrow/.style={-{Latex[length=2.3mm]}, thick}
  ]
    \node[classbox] (class) {
      \textbf{Class: Melody}\\[1mm]
      Attributes:\\
      \quad title\\
      \quad events\\
      \quad key\\[1mm]
      Methods:\\
      \quad add\_event()\\
      \quad get\_length()
    };

    \node[objectbox, right=28mm of class] (object) {
      \textbf{Object: melody\_a}\\[1mm]
      title = ``Example theme''\\
      events = (c4, d4, e4, g4, \ldots)\\
      key = C major
    };

    \draw[arrow] (class) -- node[above, align=center] {instantiated as\\a concrete melody} (object);
  \end{tikzpicture}%
  }
  \caption{A simple illustration of the difference between a class and an object. A class describes what kind of data and behaviour an object has, while an object stores one concrete instance.}
  \label{fig:oop_class_object}
\end{figure}


\subsection{Encapsulation and abstraction}

  Two central ideas in \ac{OOP} are encapsulation and abstraction. Encapsulation means that related data and behaviour are grouped together. For example, instead of storing note events in one place and melody-related operations in another, they can be kept together in a single object. This makes the program easier to reason about because the responsibilities of each object become clearer.

  Encapsulation also makes it easier for an object to preserve its own consistency. If a melody should only be changed through well-defined operations such as adding an event, transposing the line, or exporting it, then the class can ensure that the internal data remains valid. Outside code does not need to know exactly how the events are stored in order to use the melody object safely.

  Abstraction means that a class can expose the parts of the system that are important for the user of the class, while hiding unnecessary implementation details. In a well-structured program this reduces complexity, because a developer can interact with a class through a small set of meaningful operations instead of having to understand every internal detail. In a music-related program this is especially useful because it allows the code to work with concepts such as melodies, motifs, and phrases directly instead of only with low-level pitch and duration lists. Figure~\ref{fig:oop_encapsulation_abstraction} illustrates this by separating a small public interface from the internal details stored inside the object.

  \begin{figure}[htbp]
  \centering
  \resizebox{0.95\textwidth}{!}{%
  \begin{tikzpicture}[
    >=Latex,
    font=\small,
    node distance=20mm and 14mm,
    outerbox/.style={
      draw,
      rounded corners=3pt,
      minimum width=78mm,
      minimum height=50mm,
      fill=gray!6
    },
    innerbox/.style={
      draw,
      rounded corners=3pt,
      minimum width=60mm,
      align=left
    },
    caller/.style={
      draw,
      rounded corners=3pt,
      minimum width=30mm,
      align=center,
      fill=blue!6
    },
    arrow/.style={-{Latex[length=2.3mm]}, thick}
  ]
    \node[caller] (user) {Other code};

    \node[outerbox, right=24mm of user] (melody) {};
    \node[anchor=north] at (melody.north) {\textbf{Melody object}};

    \node[innerbox, fill=green!8, anchor=north] (public) at ([yshift=-6mm]melody.north) {
      \textbf{Public operations}\\
      \quad add\_event()\\
      \quad transpose()\\
      \quad to\_lilypond()
    };

    \node[innerbox, fill=yellow!10, anchor=south] (internal) at ([yshift=2mm]melody.south) {
      \textbf{Internal details}\\
      \quad events\\
      \quad cached length\\
      \quad harmony references
    };

    \draw[arrow] (user) -- node[above] {uses} (public.west);
  \end{tikzpicture}%
  }
  \caption{Illustration of encapsulation and abstraction. Encapsulation groups related data and behaviour inside one object, while abstraction lets other code use a small public interface without needing direct access to the internal representation.}
  \label{fig:oop_encapsulation_abstraction}
\end{figure}


\subsection{Composition and inheritance}

  Another important idea in \ac{OOP} is the relationship between objects. One common relationship is \textit{composition}, where one object consists of one or more other objects. A simple example is that a melody can consist of several note objects, or that a phrase can consist of one or more motifs. This is often a natural way to model structured data because the software structure can follow the same part-whole structure as the domain itself.

  A second common relationship is \textit{inheritance}, where a class derives behaviour from a more general parent class. Inheritance can be useful when several classes share a lot of common structure or should follow the same general interface. However, inheritance is not always the best solution. In many cases it is simpler and clearer to combine several smaller objects instead of building a deep hierarchy of subclasses. For this reason composition often becomes the more direct way to describe musical structure, while inheritance is more useful for shared behaviour between related classes. Figure~\ref{fig:oop_composition_inheritance} shows this contrast by placing a part-whole example beside a simple inheritance hierarchy.

  \begin{figure}[htbp]
  \centering
  \resizebox{0.85\textwidth}{!}{%
  \begin{tikzpicture}[
    >=Latex,
    font=\small,
    node distance=10mm and 16mm,
    outerbox/.style={
      draw,
      rounded corners=3pt,
      fill=gray!6
    },
    box/.style={
      draw,
      rounded corners=3pt,
      align=center,
      minimum width=30mm,
      minimum height=9mm
    },
    comp/.style={box, fill=blue!6},
    inh/.style={box, fill=green!8},
    arrow/.style={-{Latex[length=2.3mm]}, thick}
  ]
    \node[outerbox, minimum width=56mm, minimum height=54mm] (compgroup) at (-3.8,-0.4) {};
    \node[anchor=north] at (compgroup.north) {\textbf{Composition}};

    \node[comp] (phrase) at (-3.8,1.0) {Phrase};
    \node[comp] (motif) at (-3.8,-0.4) {Motif};
    \node[comp] (note) at (-3.8,-1.8) {NoteEvent};

    \draw[arrow] (phrase) -- node[right] {contains} (motif);
    \draw[arrow] (motif) -- node[right] {contains} (note);

    \node[outerbox, minimum width=76mm, minimum height=54mm] (inhgroup) at (4.2,-0.4) {};
    \node[anchor=north] at (inhgroup.north) {\textbf{Inheritance}};

    \node[inh] (base) at (4.2,0.5) {MusicalObject};
    \node[inh] (melody) at (2.5,-1.2) {Melody};
    \node[inh] (chord) at (5.9,-1.2) {ChordProgression};

    \draw[arrow] (melody) -- (base);
    \draw[arrow] (chord) -- (base);
  \end{tikzpicture}%
  }
  \caption{Comparison of two common relationships in object-oriented design. Composition models part-whole structures such as phrases, motifs, and note events, while inheritance models an is-a relationship where several classes share behaviour from a more general parent class.}
  \label{fig:oop_composition_inheritance}
\end{figure}


\subsection{Polymorphism}

  Polymorphism means that several different classes can be used through the same interface. This allows a program to treat them uniformly even though their internal behaviour differs. Figure~\ref{fig:oop_iter3_model} shows a simple example of this idea. Different constraint classes can implement the same operation, for instance checking whether a musical candidate satisfies a certain rule, while each class applies a different criterion.

  \begin{figure}[htbp]
  \centering
  \resizebox{\textwidth}{!}{%
  \begin{tikzpicture}[
    >=Latex,
    font=\small,
    node distance=10mm and 14mm,
    class/.style={
      draw,
      rounded corners=3pt,
      align=left,
      fill=gray!8,
      minimum width=42mm
    },
    support/.style={
      draw,
      rounded corners=3pt,
      align=left,
      fill=blue!6,
      minimum width=34mm
    },
    arrow/.style={-{Latex[length=2.3mm]}, thick}
  ]
    \node[class] (generator) {
      \textbf{Generator}\\
      create()\\
      evaluate()\\
      export()
    };

    \node[class, left=24mm of generator] (melody) {
      \textbf{Melody}\\
      title\\
      key\\
      events: Note[]
    };

    \node[support, below=of melody] (note) {
      \textbf{Note}\\
      pitch\\
      duration
    };

    \node[support, right=24mm of generator] (constraint) {
      \textbf{Constraint}\\
      check(candidate)
    };

    \node[support, above right=of constraint] (range) {
      \textbf{RangeConstraint}\\
      checks pitch range
    };

    \node[support, below right=of constraint] (cadence) {
      \textbf{CadenceConstraint}\\
      checks phrase ending
    };

    \draw[arrow] (generator) -- node[above] {creates} (melody);
    \draw[arrow] (note) -- node[left] {contained in} (melody);
    \draw[arrow] (constraint) -- node[above] {used by} (generator);
    \draw[arrow] (range) -- (constraint);
    \draw[arrow] (cadence) -- (constraint);
  \end{tikzpicture}%
  }
  \caption{A simple object-oriented example showing two important ideas. Composition is used when a melody consists of note objects, while polymorphism is used when different constraint classes can be treated through one shared interface.}
  \label{fig:oop_iter3_model}
\end{figure}


  This is useful because it allows a system to be extended without rewriting the code that uses the interface. New rules or behaviours can be introduced by defining a new class that follows the same pattern as the existing ones.

\subsection{Relevance to the thesis}

  In this thesis \ac{OOP} is relevant because the project works with structured musical concepts which need to be represented in an understandable manner. Without this, the code will be harder to understand, use and also to maintain. Notes, melodies, motifs, harmonic structures, and constraints are easier to reason about when they are represented as explicit objects. This does not mean that \ac{OOP} alone solves the problem of generating good music, but it provides a useful way of organising the software such that the program structure can reflect the musical structure being discussed.


\subsection{Strengths and limitations}

  The main strength of \ac{OOP} is that it can make a large program easier to structure, understand, and extend. By giving different responsibilities to different classes, the code can be decomposed into smaller and more meaningful parts. This is especially useful when the problem domain itself has many distinguishable entities.

  At the same time, \ac{OOP} does not guarantee good design. Poorly chosen classes can make the program more confusing. Too much abstraction can make the code hard to read, while unnecessary inheritance can create rigid structures that are hard to modify.
